%% parts/p08/modal-logics.tex --- Part 8.2
%% Rules remain in registry/calculi; proof material lives in annex/p08/.
\chapter{Modal Logics}\label{sec:p08-modal-logics}

\section{Language, structures, and consequence}

Modal formulas extend the classical propositional grammar by the constructor
\(\Box A\), with \(\Diamond A\defeq\neg\Box\neg A\).  These symbols acquire
no application-specific reading until a frame class and satisfaction relation
are fixed.

\begin{hbdefinition}{def:p08-modal-kripke-structure}{Modal frames and structures}
A \emph{modal Kripke frame} is a pair \(\mathfrak F=(W,R)\), where
\(W\ne\varnothing\) and \(R\subseteq W\times W\).  A \emph{modal Kripke
structure} is \(\mathcal M=(W,R,V)\), where
\(V:\mathsf{Prop}\to\mathcal P(W)\).  Satisfaction is classical for the
propositional connectives, and
\[
 \mathcal M,w\entails\Box A
 \quad\text{ iff }\quad
 \forall v\,(wRv\to\mathcal M,v\entails A).
\]
Consequently, \(\mathcal M,w\entails\Diamond A\) exactly when some
\(R\)-successor of \(w\) satisfies \(A\).  The relation \(R\) is the frame's
accessibility relation; it is not a derivability relation
\cite{Kripke1963}.
\end{hbdefinition}

\begin{hbdefinition}{def:p08-modal-validity}{Structure validity and frame validity}
A formula is \emph{valid in a structure} when it is satisfied at every world
of that structure.  It is \emph{valid on a frame} when it is valid in every
structure based on that frame.  For a class \(\mathcal C\) of frames,
\(\entails_{\mathcal C}A\) means that \(A\) is valid on every frame in
\(\mathcal C\).
\end{hbdefinition}

\begin{hbdefinition}{def:p08-modal-local-global}{Local and global modal consequence}
For a frame class \(\mathcal C\), local consequence
\(\Gamma\entails^{\ell}_{\mathcal C}A\) means that, in every structure on a
frame in \(\mathcal C\) and at every world \(w\), satisfaction of every member
of \(\Gamma\) at \(w\) implies satisfaction of \(A\) at \(w\).  Global
consequence \(\Gamma\entails^{g}_{\mathcal C}A\) means that, in every such
structure, truth of every member of \(\Gamma\) at every world implies truth of
\(A\) at every world.  The two relations can differ; every adequacy statement
below names which one it uses.
\end{hbdefinition}

\begin{hbconvention}{def:p08-modal-derivability}{Premise-sensitive modal derivability}
For any displayed normal modal sheet and any premise set \(\Gamma\), an
\(\mathfrak L\)-derivation is a finite sequence whose lines are premises,
instances of schemas on that sheet or its named base sheets, modus-ponens
consequences, or boxes of formulas having a separate no-premise derivation.
Write \(\Gamma\derives_{\mathfrak L}A\) when the sequence ends in \(A\).
Necessitation is therefore theorem-only even inside a derivation from premises;
no card below licenses premise-dependent necessitation.
\end{hbconvention}

\section{Normal modal systems}

\begin{hbwarning}{def:p08-modal-b-convention}{The system \(KTB\), the schema B, and the alias \(B\)}
Chellas writes \(KTB\) for the reflexive symmetric modal-cube system: it
contains both the T schema and the B schema.  The B schema alone corresponds
to symmetry, and Chellas writes \(KB\) for the strictly weaker system obtained
by adding only that schema to \(K\) \cite{Chellas1980}.  Hughes and Cresswell
use \(B\) as a name for the reflexive symmetric system
\cite{HughesCresswell1996}; here \(B\) is therefore recorded only as their
source-bound alias for the card whose Chellas name is \(KTB\).
\end{hbwarning}

\begin{hbdefinition}{def:p08-modal-frame-properties}{Frame properties used in the system catalog}
For a modal frame \((W,R)\), the relation \(R\) is \emph{serial} when
\(\forall x\,\exists y\,Rxy\), \emph{Euclidean} when
\(\forall x\forall y\forall z\,(Rxy\land Rxz\to Ryz)\),
\emph{directed} (or \emph{convergent}) when
\(Rxy\land Rxz\to\exists u\,(Ryu\land Rzu)\), and \emph{weakly connected}
when \(Rxy\land Rxz\to(Ryz\lor Rzy)\).  These are exact first-order
properties of the accessibility relation, not names for proof rules.  They are
needed because the catalog distinguishes \(KD\), \(KD45\), \(S4.2\), and
\(S4.3\) from systems
characterized only by reflexivity, symmetry, or transitivity; the same clauses
reappear in the correspondence table below.
\end{hbdefinition}

\begin{table}[!htbp]
\centering
\small
\begin{tabular}{@{}p{.09\textwidth}p{.18\textwidth}p{.14\textwidth}p{.27\textwidth}p{.16\textwidth}@{}}
\toprule
System & Sheet & Axioms & Frames & Parents \\
\midrule
\(K\) & \texttt{hilbert-k} & K & all frames & --- \\
\(KD\) & \texttt{hilbert-d} & K, D & serial & K \\
\(KT\) & \texttt{hilbert-t} & K, T & reflexive & KD \\
\(KTB\) & \texttt{hilbert-b} & K, T, B & reflexive, symmetric & KT \\
\(K4\) & \texttt{hilbert-k4} & K, 4 & transitive & K \\
\(S4\) (\(KT4\)) & \texttt{hilbert-s4} & K, T, 4 & reflexive, transitive & KT, K4 \\
\(KD45\) & \texttt{hilbert-kd45} & K, D, 4, 5 & serial, transitive, Euclidean & KD, K4 \\
\(S4.2\) & \texttt{hilbert-s42} & S4, .2 & directed S4 frames & S4 \\
\(S4.3\) & \texttt{hilbert-s43} & S4, .3 & weakly connected S4 frames & S4.2 \\
\(S5\) & \texttt{hilbert-s5} & K, T, 5 & equivalence frames & KTB, KD45, S4.3 \\
\(\mathsf{GL}\) & \texttt{hilbert-gl} & K, L\"ob & finite transitive irreflexive & K4 \\
\(\mathsf{Grz}\) & \texttt{hilbert-grz} & S4, Grz & finite partial orders & S4 \\
\bottomrule
\end{tabular}
\caption{The named modal systems, exact rule surfaces, and inclusion parents
used by the local lattice (apparatus A-076).  Extension names point to schemas
on the named sheets; this table does not redisplay rules.  GL and Grz are
separate anchored branches, not ordinary modal-cube nodes.}
\label{tab:p08-modal-systems}
\end{table}

\begin{systementry}{entry:p08-modal-k}{Normal modal logic \(K\)}
  \index{modal logic!K}
  \entrylanguage{Classical modal propositional formulas with \(\Box\)
    primitive and \(\Diamond A\defeq\neg\Box\neg A\).}
  \entrysemantics{All nonempty Kripke frames.}
  \entrycalculus{sheet:hilbert-k}
  \entryadequacy{cited=BlackburnDeRijkeVenema2001}{Sound and strongly complete
    for local consequence on all frames.  The inherited local sketch is kept
    separately in Theorem~\ref{thm:p08-modal-four-system-adequacy}.}
  \entrydecidability{cited=ChagrovZakharyaschev1997}{\(K\) has the finite-model
    property and a decision procedure; Section~\ref{sec:p08-modal-filtration}
    records the filtration boundary.}
  \entrycharacteristic{The weakest normal system in this displayed catalog;
    it does not validate T, 4, or 5 on arbitrary frames.}
  \entrytranslations{A compositional first-order comparison is given later,
    after its target signature and recursion have been defined.}
  \entrysources{Kripke1963,BlackburnDeRijkeVenema2001,ChagrovZakharyaschev1997}
\end{systementry}

\begin{systementry}{entry:p08-modal-d}{Normal modal logic \(KD\)}
  \index{modal logic!KD}
  \entryaliases{\(D\) in the notation of Hughes and Cresswell
    \cite{HughesCresswell1996}.}
  \entrylanguage{The modal propositional language of \(K\).}
  \entrysemantics{Serial Kripke frames: every world has an accessible world.}
  \entrycalculus{sheet:hilbert-d}
  \entryadequacy{cited=Chellas1980}{Sound and strongly complete for local
    consequence on serial frames.}
  \entrydecidability{cited=ChagrovZakharyaschev1997}{\(KD\) has the finite-model
    property and is decidable.}
  \entrycharacteristic{Its defining D schema corresponds to seriality.}
  \entrytranslations{The standard translation sends seriality to the first
    row of Table~\ref{tab:p08-modal-correspondence}.}
  \entrysources{Chellas1980,HughesCresswell1996,ChagrovZakharyaschev1997}
\end{systementry}

\begin{systementry}{entry:p08-modal-t}{Normal modal logic \(KT\)}
  \index{modal logic!KT}
  \entryaliases{\(T\) in the notation of Hughes and Cresswell
    \cite{HughesCresswell1996}.}
  \entrylanguage{The modal propositional language of \(K\).}
  \entrysemantics{Reflexive Kripke frames.}
  \entrycalculus{sheet:hilbert-t}
  \entryadequacy{cited=BlackburnDeRijkeVenema2001}{Sound and strongly complete
    for local consequence on reflexive frames.  The inherited local sketch is
    kept in Theorem~\ref{thm:p08-modal-four-system-adequacy}.}
  \entrydecidability{cited=ChagrovZakharyaschev1997}{\(KT\) has the finite-model
    property and is decidable.}
  \entrycharacteristic{Its defining T schema corresponds to reflexivity and
    entails the D schema.}
  \entrytranslations{The standard translation identifies the reflexivity
    condition in Table~\ref{tab:p08-modal-correspondence}.}
  \entrysources{BlackburnDeRijkeVenema2001,Chellas1980,ChagrovZakharyaschev1997}
\end{systementry}

\begin{systementry}{entry:p08-modal-b}{Modal-cube system \(KTB\)}
  \index{modal logic!KTB}
  \entryaliases{\(B\) in the notation of Hughes and Cresswell
    \cite{HughesCresswell1996}; not Chellas's distinct symmetry-only logic
    \(KB\).}
  \entrylanguage{The modal propositional language of \(K\).}
  \entrysemantics{Reflexive and symmetric Kripke frames.}
  \entrycalculus{sheet:hilbert-b}
  \entryadequacy{cited=HughesCresswell1996}{Sound and strongly complete for
    local consequence on reflexive symmetric frames.}
  \entrydecidability{cited=ChagrovZakharyaschev1997}{\(KTB\) has the finite-model
    property and is decidable.}
  \entrycharacteristic{The T component supplies reflexivity; the separately
    named B schema supplies symmetry.}
  \entrytranslations{The B-schema correspondence is the symmetry row of
    Table~\ref{tab:p08-modal-correspondence}.}
  \entrysources{Chellas1980,HughesCresswell1996,ChagrovZakharyaschev1997}
\end{systementry}

\begin{systementry}{entry:p08-modal-k4}{Normal modal logic \(K4\)}
  \index{modal logic!K4}
  \entrylanguage{The modal propositional language of \(K\).}
  \entrysemantics{Transitive Kripke frames.}
  \entrycalculus{sheet:hilbert-k4}
  \entryadequacy{cited=Chellas1980}{Sound and strongly complete for local
    consequence on transitive frames.}
  \entrydecidability{cited=ChagrovZakharyaschev1997}{\(K4\) has the finite-model
    property and is decidable.}
  \entrycharacteristic{Its defining 4 schema corresponds to transitivity.}
  \entrytranslations{The standard translation identifies transitivity in
    Table~\ref{tab:p08-modal-correspondence}.}
  \entrysources{Chellas1980,HughesCresswell1996,ChagrovZakharyaschev1997}
\end{systementry}

\begin{systementry}{entry:p08-modal-s4}{Normal modal logic \(S4\) (\(KT4\))}
  \index{modal logic!S4}
  \entrylanguage{The modal propositional language of \(K\).}
  \entrysemantics{Reflexive and transitive Kripke frames.}
  \entrycalculus{sheet:hilbert-s4}
  \entryadequacy{cited=BlackburnDeRijkeVenema2001}{Sound and strongly complete
    for local consequence on reflexive transitive frames.  The inherited local
    sketch is kept in Theorem~\ref{thm:p08-modal-four-system-adequacy}.}
  \entrydecidability{cited=ChagrovZakharyaschev1997}{\(S4\) has the finite-model
    property and is decidable.}
  \entrycharacteristic{It combines the frame conditions represented by T and
    4; neither the .2 nor .3 extension is built into \(S4\).}
  \entrytranslations{The G\"odel--McKinsey--Tarski embedding of IPL is recorded
    in Theorem~\ref{thm:p08-ipl-s4-embedding}.}
  \entrysources{BlackburnDeRijkeVenema2001,HughesCresswell1996,ChagrovZakharyaschev1997}
\end{systementry}

\begin{systementry}{entry:p08-modal-kd45}{Normal modal logic \(KD45\)}
  \index{modal logic!KD45}
  \entrylanguage{The modal propositional language of \(K\).}
  \entrysemantics{Serial, transitive, and Euclidean Kripke frames.}
  \entrycalculus{sheet:hilbert-kd45}
  \entryadequacy{cited=ChagrovZakharyaschev1997}{Sound and strongly complete
    for local consequence on serial transitive Euclidean frames.}
  \entrydecidability{cited=ChagrovZakharyaschev1997}{\(KD45\) has the finite-model
    property and is decidable.}
  \entrycharacteristic{Its D, 4, and 5 components correspond respectively to
    seriality, transitivity, and Euclideanness.}
  \entrytranslations{Each component has the standard first-order
    correspondence in Table~\ref{tab:p08-modal-correspondence}.}
  \entrysources{Chellas1980,HughesCresswell1996,ChagrovZakharyaschev1997}
\end{systementry}

\begin{systementry}{entry:p08-modal-s42}{Normal modal logic \(S4.2\)}
  \index{modal logic!S4.2}
  \entrylanguage{The modal propositional language of \(K\).}
  \entrysemantics{Reflexive transitive frames directed above each world: any
    two successors have a common successor.}
  \entrycalculus{sheet:hilbert-s42}
  \entryadequacy{cited=ChagrovZakharyaschev1997}{Sound and strongly complete
    for local consequence on the stated directed \(S4\) frames.}
  \entrydecidability{cited=ChagrovZakharyaschev1997}{\(S4.2\) has the finite-model
    property and is decidable.}
  \entrycharacteristic{The .2 schema adds directed convergence to \(S4\).}
  \entrytranslations{Its exact relational condition is the .2 row of
    Table~\ref{tab:p08-modal-correspondence}.}
  \entrysources{LemmonScott1977,HughesCresswell1996,ChagrovZakharyaschev1997}
\end{systementry}

\begin{systementry}{entry:p08-modal-s43}{Normal modal logic \(S4.3\)}
  \index{modal logic!S4.3}
  \entrylanguage{The modal propositional language of \(K\).}
  \entrysemantics{Reflexive transitive frames whose successors at each world
    are weakly connected.}
  \entrycalculus{sheet:hilbert-s43}
  \entryadequacy{cited=ChagrovZakharyaschev1997}{Sound and strongly complete
    for local consequence on the stated weakly connected \(S4\) frames.}
  \entrydecidability{cited=ChagrovZakharyaschev1997}{\(S4.3\) has the finite-model
    property and is decidable.}
  \entrycharacteristic{The .3 schema adds weak connectedness to \(S4\) and
    contains \(S4.2\) in the displayed selection.}
  \entrytranslations{Its exact relational condition is the .3 row of
    Table~\ref{tab:p08-modal-correspondence}.}
  \entrysources{LemmonScott1977,HughesCresswell1996,ChagrovZakharyaschev1997}
\end{systementry}

\begin{systementry}{entry:p08-modal-s5}{Normal modal logic \(S5\)}
  \index{modal logic!S5}
  \entrylanguage{The modal propositional language of \(K\).}
  \entrysemantics{Kripke frames whose accessibility relation is an equivalence
    relation.}
  \entrycalculus{sheet:hilbert-s5}
  \entryadequacy{cited=BlackburnDeRijkeVenema2001}{Sound and strongly complete
    for local consequence on equivalence frames.  The inherited local sketch
    is kept in Theorem~\ref{thm:p08-modal-four-system-adequacy}.}
  \entrydecidability{cited=ChagrovZakharyaschev1997}{\(S5\) has the finite-model
    property and is decidable.}
  \entrycharacteristic{The registered sheet is T plus 5; reflexivity together
    with Euclideanness yields symmetry and transitivity.}
  \entrytranslations{The T and 5 standard translations give the corresponding
    frame conditions in Table~\ref{tab:p08-modal-correspondence}.}
  \entrysources{BlackburnDeRijkeVenema2001,Chellas1980,ChagrovZakharyaschev1997}
\end{systementry}

\begin{systementry}{entry:p08-modal-gl}{G\"odel--L\"ob provability logic GL}
  \index{provability logic!GL}
  \entryaliases{GL; a provability logic, displayed outside the ordinary modal
    cube.}
  \entrylanguage{The modal propositional language of \(K\).}
  \entrysemantics{For Kripke theoremhood completeness: finite transitive
    irreflexive frames.  No arbitrary-frame strong-completeness claim is made.}
  \entrycalculus{sheet:hilbert-gl}
  \entryadequacy{cited=ChagrovZakharyaschev1997}{Theoremhood is sound and
    complete for finite transitive irreflexive frames; the ordinary canonical
    construction below is not asserted for GL.}
  \entrydecidability{cited=Boolos1993}{GL has the finite-model property and is
    decidable.}
  \entrycharacteristic{L\"ob's schema; the arithmetical interpretation is not
    an added edge in the modal cube.}
  \entrytranslations{See Solovay completeness,
    Theorem~\ref{thm:p08-gl-arithmetical-completeness}.}
  \entrysources{ChagrovZakharyaschev1997,Solovay1976,Boolos1993}
\end{systementry}

\begin{systementry}{entry:p08-modal-grz}{Grzegorczyk modal logic Grz}
  \index{modal logic!Grzegorczyk logic}
  \entryaliases{Grz; displayed as a separate extension anchored at \(S4\).}
  \entrylanguage{The modal propositional language of \(K\).}
  \entrysemantics{Finite partial-order frames for theoremhood completeness;
    accessibility is reflexive, transitive, and antisymmetric.}
  \entrycalculus{sheet:hilbert-grz}
  \entryadequacy{cited=ChagrovZakharyaschev1997}{Cited theoremhood completeness
    for finite partial orders; the ordinary canonical method is not asserted.}
  \entrydecidability{cited=ChagrovZakharyaschev1997}{Grz has the finite-model
    property and is decidable.}
  \entrycharacteristic{Its Grz schema extends \(S4\).  No inclusion edge from
    Grz to \(S4.2\), \(S4.3\), or \(S5\) is asserted here.}
  \entrytranslations{The relational completeness statement is kept separate
    from the ordinary canonical-model method.}
  \entrysources{LemmonScott1977,ChagrovZakharyaschev1997}
\end{systementry}

\begin{hbmethod}{def:p08-modal-canonical-model}{Canonical-model method}
\index{canonical-model method!modal logic}
For a fixed normal modal logic \(\mathfrak L\), let the worlds be the maximal
\(\mathfrak L\)-consistent sets.  Define the metalinguistic proof relation
\(R_{\mathfrak L}\) by
\[
 \Gamma R_{\mathfrak L}\Delta
 \quad\text{ iff }\quad
 \forall A\,(\Box A\in\Gamma\to A\in\Delta),
\]
and make an atom true at \(\Gamma\) exactly when it belongs to \(\Gamma\).
This construction is needed to distinguish the proof model from an arbitrary
semantic structure.  Its worlds are theories, unlike filtration worlds, which
are equivalence classes of worlds from a given model.  The relation is used by
the modal existence lemma, truth lemma, and the frame-property checks in the
following adequacy sketch.  The ordinary construction applies to the canonical
systems \(K\), \(KD\), \(KT\), \(KTB\), \(K4\), \(S4\) (\(KT4\)),
\(KD45\), \(S4.2\), \(S4.3\), and \(S5\) with their source-defined
schemas; it is not the completeness route asserted for GL or Grz.
\end{hbmethod}

\begin{hbtheorem}{thm:p08-modal-four-system-adequacy}{sketched=annex:p8.modal.1}{Soundness and strong completeness for \(K\), \(KT\), \(S4\), and \(S5\)}
For \(\mathfrak L\in\{K,KT,S4,S5\}\), let
\(\mathcal C_{\mathfrak L}\) be respectively all frames, reflexive frames,
reflexive transitive frames, and equivalence frames.  For every premise set
\(\Gamma\) and modal formula \(A\),
\[
 \Gamma\derives_{\mathfrak L}A
 \quad\text{ iff }\quad
 \Gamma\entails^{\ell}_{\mathcal C_{\mathfrak L}}A.
\]
Thus the statement is about strong completeness for local consequence.  It is
not a global-consequence theorem and its inherited source argument is a sketch.
\end{hbtheorem}

\begin{figure}[!htbp]
\centering
\begin{tikzpicture}[
  >=Latex,
  every node/.style={draw,rounded corners,inner sep=4pt,font=\small},
  x=1.8cm,y=1.05cm]
  \node (k) at (0,0) {\(K\)};
  \node (d) at (-2,1.5) {\(KD\)};
  \node (k4) at (2,1.5) {\(K4\)};
  \node (t) at (-2,3) {\(KT\)};
  \node (kd45) at (0,2.2) {\(KD45\)};
  \node (s4) at (2,3) {\(S4\)};
  \node (gl) at (4,3) {\(\mathsf{GL}\)};
  \node (b) at (-2,4.5) {\(KTB\)};
  \node (s42) at (2,4.5) {\(S4.2\)};
  \node (grz) at (4,4.5) {\(\mathsf{Grz}\)};
  \node (s43) at (2,6) {\(S4.3\)};
  \node (s5) at (0,7.5) {\(S5\)};
  \foreach \x/\y in {
    k/d,k/k4,d/t,d/kd45,k4/s4,k4/kd45,k4/gl,
    t/b,t/s4,s4/s42,s4/grz,s42/s43,s43/s5,b/s5}
    \draw[->] (\x) -- (\y);
  \draw[->,preaction={draw=white,line width=4pt}]
    (kd45) -- (s5);
\end{tikzpicture}
\caption{Theorem-set inclusion among only the twelve displayed systems
(apparatus A-136).  This is the Hasse diagram for the displayed inclusions in
Table~\ref{tab:p08-modal-systems}, not a complete lattice of normal modal
logics.  GL and Grz are separately anchored at \(K4\) and \(S4\); no unshown
comparability with \(S5\), \(S4.2\), or \(S4.3\) is asserted.}
\label{fig:p08-modal-lattice}
\end{figure}

\section{Correspondence and translations}

\begin{hbdefinition}{def:p08-modal-standard-translation}{First-order standard translation}
For each propositional variable \(p\), choose a unary predicate symbol
\(P_p\).  The standard translation at the displayed free world variable
\(x\), written \(\operatorname{ST}_x\), is the metalinguistic recursion
\[
\begin{aligned}
 \operatorname{ST}_x(p)&=P_p(x),
 &\operatorname{ST}_x(\bottom)&=\bottom,\\
 \operatorname{ST}_x(\neg A)&=\neg\operatorname{ST}_x(A),
 &\operatorname{ST}_x(A\circ B)&=
   \operatorname{ST}_x(A)\circ\operatorname{ST}_x(B),\\
 \operatorname{ST}_x(\Box A)&=
   \forall y\,(Rxy\to\operatorname{ST}_y(A)),
 &\operatorname{ST}_x(\Diamond A)&=
   \exists y\,(Rxy\land\operatorname{ST}_y(A)),
\end{aligned}
\]
where \(\circ\in\{\land,\lor,\to\}\) and \(y\) is fresh.  This is a
translation into first-order formulas over the frame relation and valuation
predicates, not a new modal connective.  It is needed to state the exact frame
conditions below and later to formulate the bisimulation characterization.
\end{hbdefinition}

\begin{hbtheorem}{thm:p08-modal-correspondence}{cited=Chellas1980}{Elementary modal frame correspondences}
For every Kripke frame, validity of each named schema in the first column of
Table~\ref{tab:p08-modal-correspondence} is equivalent to the displayed
first-order condition on its accessibility relation.  The row named B concerns
the B schema alone; the displayed system \(KTB\) additionally contains T and
hence uses reflexive symmetric frames.
\end{hbtheorem}

\begin{table}[!htbp]
\centering
\small
\begin{tabular}{@{}lll@{}}
\toprule
Schema & First-order condition on \(R\) & Conventional name \\
\midrule
D & \(\forall x\,\exists y\,Rxy\) & serial \\
T & \(\forall x\,Rxx\) & reflexive \\
B & \(\forall x\forall y\,(Rxy\to Ryx)\) & symmetric \\
4 & \(\forall x\forall y\forall z\,(Rxy\land Ryz\to Rxz)\) & transitive \\
5 & \(\forall x\forall y\forall z\,(Rxy\land Rxz\to Ryz)\) & Euclidean \\
.2 & \(Rxy\land Rxz\to\exists u\,(Ryu\land Rzu)\) & directed/convergent \\
.3 & \(Rxy\land Rxz\to(Ryz\lor Rzy)\) & weakly connected \\
\bottomrule
\end{tabular}
\caption{Local frame-correspondence data for D-P10-003 and N-b-002.  The .3
row is used on the transitive S4 presentation.  This Phase-4 table supplies
local data only; it does not close the Phase-6 master table A-077 or its
algebra column.}
\label{tab:p08-modal-correspondence}
\end{table}

\begin{hbdefinition}{def:p08-sahlqvist-class}{The Sahlqvist class used here}
Rewrite modal formulas using \(\bottom,\topv,\neg,\land,\lor,\Box,\Diamond\).
An occurrence of a variable is positive when it lies under an even number of
negations and negative when it lies under an odd number.  A \emph{boxed atom}
is \(\Box^n p\) for \(n\geq0\).  A \emph{Sahlqvist antecedent} is generated
from boxed atoms, formulas all of whose variable occurrences are negative,
and \(\bottom,\topv\) by \(\land,\lor,\Diamond\).  A \emph{Sahlqvist
implication} has such an antecedent and a consequent all of whose variable
occurrences are positive; finite conjunctions of these implications form the
Sahlqvist class fixed here.  This exact syntactic class is narrower than the
class of all modal formulas and is introduced because the following theorem
does not hold for arbitrary schemas.
\end{hbdefinition}

\begin{hbtheorem}{thm:p08-sahlqvist}{cited=Sahlqvist1975}{Sahlqvist theorem}
Every formula in the class fixed by
Definition~\ref{def:p08-sahlqvist-class} has an effectively computable
first-order frame correspondent, is canonical, and axiomatizes a normal modal
logic complete for the corresponding elementary frame class.  No claim is
made here that an arbitrary modal axiom is a Sahlqvist formula.
\end{hbtheorem}

\begin{hbdefinition}{def:p08-ipl-s4-translation}{G\"odel--McKinsey--Tarski translation}
For an IPL formula \(A\), define its modal translation \(A^{\Box}\) by
\(p^{\Box}\defeq\Box p\), \(\bottom^{\Box}\defeq\bottom\), commuting with
\(\land\) and \(\lor\), and
\[
 (A\to B)^{\Box}\defeq\Box(A^{\Box}\to B^{\Box}).
\]
This metalinguistic recursion is distinct from both the first-order standard
translation and the G\"odel--Gentzen negative translation of Chapter~8.1.
\end{hbdefinition}

\begin{hbtheorem}{thm:p08-ipl-s4-embedding}{cited=Godel1933}{Embedding IPL into \(S4\)}
For every IPL formula \(A\),
\(\derives_{\mathsf{IPL}}A\) exactly when
\(\derives_{S4}A^{\Box}\).
\end{hbtheorem}

\begin{hbtheorem}{thm:p08-gl-arithmetical-completeness}{cited=Solovay1976}{Arithmetical completeness of GL}
An arithmetical realization maps propositional variables to closed
arithmetical sentences, commutes with the Boolean connectives, and interprets
\(\Box B\) by the standard provability predicate
\(\mathrm{Prov}_{\mathsf{PA}}\) applied to the G\"odel number of the realized
sentence \(B\).  A modal formula is derivable in GL exactly when its image
under every such arithmetical realization is derivable in \(\mathsf{PA}\).  This is an
arithmetical-interpretation theorem, not the finite-frame completeness claim on
the GL card \cite{Boolos1993}.
\end{hbtheorem}

\begin{table}[!htbp]
\centering
\small
\begin{tabular}{@{}p{.17\textwidth}p{.14\textwidth}p{.24\textwidth}p{.34\textwidth}@{}}
\toprule
Translation & Source & Target & Exact role \\
\midrule
\(\operatorname{ST}_x\) & basic modal logic & first-order frame language &
correspondence and van Benthem theorem \\
\((\cdot)^{\Box}\) & IPL & \(S4\) & theoremhood embedding \\
arithmetical realization & GL & arithmetic with \(\mathrm{Prov}_{\mathsf{PA}}\) &
Solovay completeness \\
\bottomrule
\end{tabular}
\caption{Three noninterchangeable modal translations (apparatus A-083).}
\label{tab:p08-modal-translations}
\end{table}

\section{Bisimulation}

\begin{hbdefinition}{def:p08-modal-bisimulation}{Bisimulation of pointed Kripke structures}
Let \(\mathcal M=(W,R,V)\) and
\(\mathcal N=(W',R',V')\) use the same propositional signature.  A relation
\(Z\subseteq W\times W'\) is a \emph{bisimulation} when, whenever \(wZw'\):
\begin{enumerate}
\item for every atom \(p\), \(w\in V(p)\) exactly when \(w'\in V'(p)\)
      (atoms);
\item if \(wRu\), some \(u'\) satisfies \(w'R'u'\) and \(uZu'\) (forth);
\item if \(w'R'u'\), some \(u\) satisfies \(wRu\) and \(uZu'\) (back).
\end{enumerate}
Two pointed structures \((\mathcal M,w)\) and \((\mathcal N,w')\) are
\emph{bisimilar} when some such \(Z\) relates \(w\) to \(w'\).  This relation
is needed because ordinary isomorphism is stronger than the modal
indistinguishability used by the following results.
\end{hbdefinition}

\begin{hbtheorem}{thm:p08-modal-bisimulation-invariance}{cited=BlackburnDeRijkeVenema2001}{Bisimulation invariance}
If \(wZw'\) for a bisimulation between \(\mathcal M\) and \(\mathcal N\),
then for every basic modal formula \(A\),
\[
 \mathcal M,w\entails A
 \quad\text{ iff }\quad
 \mathcal N,w'\entails A.
\]
\end{hbtheorem}

\begin{hbtheorem}{thm:p08-van-benthem}{cited=VanBenthem1976}{Van Benthem characterization theorem}
Over all Kripke structures, a first-order formula with one free world variable
is equivalent to \(\operatorname{ST}_x(A)\) for some basic modal formula
\(A\) exactly when it is invariant under bisimulations.  This statement is not
restricted to finite structures; the finite-model version is a different
theorem.
\end{hbtheorem}

\begin{figure}[!htbp]
\centering
\begin{tikzpicture}[>=Latex,node distance=2.5cm,font=\small]
  \node[draw,rounded corners] (w) {\(w\in\mathcal M\)};
  \node[draw,rounded corners,below=of w] (u) {\(u\in\mathcal M\)};
  \node[draw,rounded corners,right=4.3cm of w] (wp) {\(w'\in\mathcal N\)};
  \node[draw,rounded corners,below=of wp] (up) {\(u'\in\mathcal N\)};
  \draw[->] (w) -- node[left]{\(R\)} (u);
  \draw[->] (wp) -- node[right]{\(R'\)} (up);
  \draw[dashed,<->] (w) -- node[above]{\(Z\): atoms} (wp);
  \draw[dashed,<->] (u) -- node[below]{\(Z\): forth/back witness} (up);
\end{tikzpicture}
\caption{A bisimulation square (apparatus A-123).  Atom agreement is checked
at each related pair; forth matches every left successor, and back matches
every right successor.  The diagram depicts one matched pair, not a
functionality requirement on \(Z\).}
\label{fig:p08-modal-bisimulation}
\end{figure}

\section{Filtration and finite models}\label{sec:p08-modal-filtration}

\begin{hbmethod}{def:p08-modal-filtration}{Filtration through a finite formula set}
\index{filtration!modal logic}
Fix a finite subformula-closed set \(\Phi\) and a Kripke structure
\(\mathcal M=(W,R,V)\).  Put
\[
 w\sim_{\Phi}u
 \quad\text{ iff }\quad
 \forall A\in\Phi\,
   (\mathcal M,w\entails A\text{ iff }\mathcal M,u\entails A).
\]
The quotient has worlds \([w]_{\Phi}\), at most \(2^{\lvert\Phi\rvert}\)
of them, and the induced atomic valuation.  The minimal quotient relation
relates \([w]_{\Phi}\) to \([u]_{\Phi}\) when some representatives are
\(R\)-related.  The maximal quotient relation relates them when, for every
\(\Box A\in\Phi\), truth of \(\Box A\) at \(w\) implies truth of \(A\) at
\(u\).  A \emph{filtration relation} is a quotient relation between those two
bounds.  This construction is an image of a given semantic model, unlike the
maximal-consistent-theory construction in
Method~\ref{def:p08-modal-canonical-model}; it is used to preserve truth on
\(\Phi\) while making the model finite.
\end{hbmethod}

\begin{hbtheorem}{thm:p08-modal-filtration}{cited=Segerberg1971}{Filtration and decidability}
Every filtration as in Method~\ref{def:p08-modal-filtration} preserves the
truth of each formula in \(\Phi\).  For \(K\) use an ordinary filtration
relation; for \(KT\) use its reflexive closure; for \(S4\) use the standard
reflexive-transitive filtration construction; and for \(S5\) use the
equivalence-class construction.
In each case the declared filtration relation remains between the minimal and maximal
bounds and preserves the target frame class.  Hence \(K\), \(KT\), \(S4\),
and \(S5\) have the
finite-model property and are decidable
\cite{BlackburnDeRijkeVenema2001,ChagrovZakharyaschev1997}.  This is cited new
content: no local proof or inherited compression denominator is claimed.
\end{hbtheorem}

\section{Minimal failures of T, 4, and 5}

The source's degenerate T countermodel has one world, no accessibility edge,
and \(p\) false: \(\Box p\) is vacuously true while \(p\) is false.  The
first panel below uses a nondegenerate variant with the same failure so that
the relevant successor valuation is visible.  The other two panels retain the
source frames and valuations.

\begin{figure}[!htbp]
\centering
\begin{tikzpicture}[>=Latex,font=\scriptsize,x=.9cm,y=1.1cm,
  scale=.9,transform shape]
  \node[font=\small,draw=none] at (0,2.3) {A-139: T};
  \node[draw,circle] (tw) at (-.7,1) {\(w:\neg p\)};
  \node[draw,circle] (tu) at (.7,1) {\(u:p\)};
  \draw[->] (tw) -- node[above]{\(R\)} (tu);
  \node[draw=none,align=center] at (0,0)
    {\(w\entails\Box p\)\\\(w\not\entails p\)};

  \node[font=\small,draw=none] at (4.2,2.3) {A-140: 4};
  \node[draw,circle] (fw) at (2.5,1) {\(w\)};
  \node[draw,circle] (fu) at (4.2,1) {\(u:p\)};
  \node[draw,circle] (fv) at (5.9,1) {\(v:\neg p\)};
  \draw[->] (fw) -- (fu);
  \draw[->] (fu) -- (fv);
  \node[draw=none,align=center] at (4.2,0)
    {\(w\entails\Box p\)\\\(w\not\entails\Box\Box p\); no \(wRv\)};

  \node[font=\small,draw=none] at (9,2.3) {A-141: 5};
  \node[draw,circle] (vw) at (8,1) {\(w:p\)};
  \node[draw,circle] (vu) at (10,1) {\(u:\neg p\)};
  \draw[->] (vw) -- (vu);
  \draw[->] (vw) edge[loop left] (vw);
  \draw[->] (vu) edge[loop right] (vu);
  \node[draw=none,align=center] at (9,-.5)
    {\(w\entails\Diamond p\)\\\(w\not\entails\Box\Diamond p\); no \(uRw\)};
\end{tikzpicture}
\caption{Fully specified countermodels for the T, 4, and 5 schemas.  Only the
displayed edges are present.  The first root is not reflexive, the middle frame
is not transitive, and the last is the inherited two-world reflexive but
non-Euclidean frame.}
\label{fig:p08-modal-countermodels}
\end{figure}

\section{Neighborhood semantics}

The set-of-propositions semantics used here originates independently with
Scott and Montague \cite{Scott1970,Montague1970}.  Pacuit supplies the modern
neighborhood presentation and comparison with relational semantics
\cite{Pacuit2017}; the exact \(E\)/\(EM\)/\(EC\)/\(EN\)/\(EMC\) naming
convention and the cited
calculus results below follow Chellas \cite{Chellas1980}.

\begin{hbdefinition}{def:p08-neighborhood-model}{Neighborhood model}
A \emph{neighborhood model} is
\(\mathcal M=(W,N,V)\), where \(W\ne\varnothing\),
\(V:\mathsf{Prop}\to\mathcal P(W)\), and
\(N:W\to\mathcal P(\mathcal P(W))\).  For a formula \(A\), define
its truth-set abbreviation by
\(\lbrack\!\lbrack A\rbrack\!\rbrack_{\mathcal M}
\defeq\{w\in W:\mathcal M,w\entails A\}\).  The modal clause is
\[
 \mathcal M,w\entails\Box A
 \quad\text{ iff }\quad
 \lbrack\!\lbrack A\rbrack\!\rbrack_{\mathcal M}\in N(w).
\]
The function \(N\) is needed because a single accessibility relation
cannot represent every congruential non-normal logic.  Its standing is a
semantic assignment of proposition sets to each world; it is neither the
system \(EN\) nor the set of normal worlds used in the later
Routley--Meyer structure.  It differs from a Kripke
relation by taking subsets of \(W\), rather than worlds, as its local semantic
arguments, and is used in all five cards below
\cite{Scott1970,Montague1970,Chellas1980,Pacuit2017}.
\end{hbdefinition}

\begin{hbdefinition}{def:p08-neighborhood-conditions}{Neighborhood closure conditions}
For every world \(w\) and \(X,Y\subseteq W\), call \(N\)
\emph{upward closed} when
\(X\in N(w)\) and \(X\subseteq Y\) imply
\(Y\in N(w)\); \emph{binary-intersection closed} when
\(X,Y\in N(w)\) implies
\(X\cap Y\in N(w)\); and \emph{unit containing} when
\(W\in N(w)\).  These are semantic frame conditions, not additional
object-language rules.  They are introduced separately because Chellas's
schemas M, C, and N determine different combinations of them
\cite{Chellas1980,Pacuit2017}.
\end{hbdefinition}

\begin{table}[!htbp]
\centering
\small
\begin{tabular}{@{}p{.10\textwidth}p{.19\textwidth}p{.18\textwidth}p{.42\textwidth}@{}}
\toprule
System & Defining composition & Exact sheet & Neighborhood condition \\
\midrule
\(E\) & E & \texttt{hilbert-e} & arbitrary \\
\(EM\) & E+M & \texttt{hilbert-m} & upward closed \\
\(EC\) & E+C & \texttt{hilbert-c} & binary-intersection closed \\
\(EN\) & E+N & \texttt{hilbert-n} & contains \(W\) \\
\(EMC\) & E+M+C & \texttt{hilbert-r} & upward and intersection closed \\
\(K=EMCN\) & E+M+C+N & \texttt{hilbert-k} & upward, intersection closed, contains \(W\) \\
\bottomrule
\end{tabular}
\caption{Chellas's compositional names for the displayed non-normal systems.
M, C, and N in the second column are axiom names, not substitute system
identifiers.  The minimal regular logic is \(EMC\), while normal \(K\) is
\(EMCN\).}
\label{tab:p08-neighborhood-systems}
\end{table}

\begin{systementry}{entry:p08-neighborhood-e}{Congruential modal logic \(E\)}
  \index{neighborhood logic!E}
  \entrylanguage{Classical modal propositional formulas with \(\Box\) and the
    defined dual \(\Diamond\).}
  \entrysemantics{Arbitrary neighborhood models as in
    Definition~\ref{def:p08-neighborhood-model}.}
  \entrycalculus{sheet:hilbert-e}
  \entryadequacy{cited=Chellas1980}{Theoremhood is complete for arbitrary
    neighborhood frames.  No strong-consequence claim is made.}
  \entrydecidability{cited=Chellas1980}{E is determined by finite
    neighborhood models and is decidable; no sharper complexity claim is made.}
  \entrycharacteristic{Classical reasoning and theorem-only replacement of
    equivalents, without M, C, or N.}
  \entrytranslations{Truth sets replace relational translation; no ordinary
    Kripke frame class is claimed adequate for E.}
  \entrysources{Scott1970,Montague1970,Chellas1980,Pacuit2017}
\end{systementry}

\begin{systementry}{entry:p08-neighborhood-m}{Monotone modal logic \(EM\)}
  \index{neighborhood logic!EM}
  \entryaliases{E plus Chellas's schema M.}
  \entrylanguage{The modal propositional language of E.}
  \entrysemantics{Neighborhood models whose neighborhood families are upward
    closed.}
  \entrycalculus{sheet:hilbert-m}
  \entryadequacy{cited=Chellas1980}{Complete at theoremhood for upward-closed
    frames; no strong-consequence claim is made.}
  \entrydecidability{cited=Chellas1980}{EM is determined by finite
    neighborhood models and is decidable; no sharper complexity claim is made.}
  \entrycharacteristic{The M extension corresponds exactly to upward closure;
    intersection closure and the unit condition are not assumed.}
  \entrytranslations{Its semantic condition is the M row of
    Table~\ref{tab:p08-neighborhood-systems}.}
  \entrysources{Chellas1980,Pacuit2017}
\end{systementry}

\begin{systementry}{entry:p08-neighborhood-c}{Intersection modal logic \(EC\)}
  \index{neighborhood logic!EC}
  \entryaliases{E plus Chellas's schema C.}
  \entrylanguage{The modal propositional language of E.}
  \entrysemantics{Neighborhood models closed under binary intersections.}
  \entrycalculus{sheet:hilbert-c}
  \entryadequacy{cited=Chellas1980}{Theoremhood is complete for
    intersection-closed neighborhood frames.  No strong-consequence claim is
    made.}
  \entrydecidability{cited=Chellas1980}{EC is determined by finite
    neighborhood models and is decidable; no sharper complexity claim is made.}
  \entrycharacteristic{The C extension corresponds exactly to binary
    intersection closure; upward closure and the unit condition are not assumed.}
  \entrytranslations{Its semantic condition is the C row of
    Table~\ref{tab:p08-neighborhood-systems}.}
  \entrysources{Chellas1980,Pacuit2017}
\end{systementry}

\begin{systementry}{entry:p08-neighborhood-n}{Unit modal logic \(EN\)}
  \index{neighborhood logic!EN}
  \entryaliases{E plus Chellas's schema N; not normal modal logic K.}
  \entrylanguage{The modal propositional language of E.}
  \entrysemantics{Neighborhood models in which every neighborhood family
    contains the full world set \(W\).}
  \entrycalculus{sheet:hilbert-n}
  \entryadequacy{cited=Chellas1980}{Cited theoremhood completeness for
    unit-containing neighborhood frames; no strong consequence claim is made.}
  \entrydecidability{cited=Chellas1980}{EN is determined by finite
    neighborhood models and is decidable; no sharper complexity claim is made.}
  \entrycharacteristic{The N extension corresponds exactly to the unit
    condition; upward and intersection closure are not assumed.}
  \entrytranslations{Its semantic condition is the N row of
    Table~\ref{tab:p08-neighborhood-systems}.}
  \entrysources{Chellas1980,Pacuit2017}
\end{systementry}

\begin{systementry}{entry:p08-neighborhood-r}{Regular modal logic \(EMC\)}
  \index{neighborhood logic!EMC}
  \entryaliases{E plus Chellas's schemas M and C.}
  \entrylanguage{The modal propositional language of E.}
  \entrysemantics{Neighborhood models that are both upward closed and closed
    under binary intersections.}
  \entrycalculus{sheet:hilbert-r}
  \entryadequacy{cited=Chellas1980}{Cited theoremhood completeness for upward-
    and intersection-closed frames; no strong consequence claim is made.}
  \entrydecidability{cited=Chellas1980}{EMC is determined by finite
    neighborhood models and is decidable; no sharper complexity claim is made.}
  \entrycharacteristic{EMC contains M and C but not N; adding N yields normal K.}
  \entrytranslations{The system map is Table~\ref{tab:p08-neighborhood-systems}.}
  \entrysources{Chellas1980,Pacuit2017}
\end{systementry}

\section{Frame and two-index boundaries}

\begin{hbtheorem}{thm:p08-kripke-incomplete-logics}{cited=Thomason1974}{Existence of Kripke-incomplete normal modal logics}
Thomason exhibited a normal modal logic that is not complete with respect to
any class of ordinary Kripke frames \cite{Thomason1974}.  Fine exhibited such
an incomplete normal extension containing S4 \cite{Fine1974}.  These two
existence results do not imply that every noncanonical logic is Kripke
incomplete; the broader boundary is surveyed in
\cite{ChagrovZakharyaschev1997}.
\end{hbtheorem}

\begin{hbdefinition}{def:p08-general-frame}{General frame boundary}
A \emph{general frame} is a triple \((W,R,\mathcal A)\), where
\((W,R)\) is a Kripke frame and \(\mathcal A\subseteq\mathcal P(W)\) is a
specified family of admissible proposition sets closed under the Boolean and
modal operations used by the language.  Validity restricts valuations to
members of \(\mathcal A\).  This extra admissible-set component distinguishes a
general frame from the ordinary frame of
Definition~\ref{def:p08-modal-kripke-structure}; it is used to explain how
relational methods can represent logics missed by unrestricted Kripke
valuations.  No blanket claim that every noncanonical logic is repaired by one
fixed general-frame class is made.
The admissible-set semantics is historically connected with Scott's framework
\cite{Scott1970}; this definition does not identify general frames with the
neighborhood models below.
\end{hbdefinition}

\begin{hbdefinition}{def:p08-two-index-actuality}{Davies--Humberstone's source operators \(A\) and \(\mathcal F\)}
Davies and Humberstone's two-index presentation of \(S5\) uses a
model \(\langle W,V\rangle\) and writes truth as
\(\langle W,V\rangle\models_x^y\alpha\).  The lower index \(x\) is the
world of evaluation and the upper index \(y\) is the actual-world
perspective.  For atoms and the three modal operators used here, their clauses
are
\[
\begin{aligned}
 \langle W,V\rangle\models_x^y p
 &\quad\text{ iff }\quad V(p,x)=1,\\
 \langle W,V\rangle\models_x^y\Box\alpha
 &\quad\text{ iff }\quad
   \forall z\in W\;\langle W,V\rangle\models_z^y\alpha,\\
 \langle W,V\rangle\models_x^y A\alpha
 &\quad\text{ iff }\quad
   \langle W,V\rangle\models_y^y\alpha,\\
 \langle W,V\rangle\models_x^y\mathcal F\alpha
 &\quad\text{ iff }\quad
   \forall z\in W\;\langle W,V\rangle\models_x^z\alpha.
\end{aligned}
\]
Boolean clauses are evaluated at the displayed indices.  Here \(A\) is the
source actuality operator and \(\mathcal F\) is the source fixedly operator;
within this definition they are unary object-language operators, while
\(\alpha\) is the source formula metavariable.  The extra index is necessary
because one index cannot record both the evaluation world and the actual-world
perspective.  The operators are nonredundant with \(\Box\): the S5 box varies
the lower evaluation index, \(A\) resets it to the upper actual-world index,
and \(\mathcal F\) instead varies the upper index while holding the lower one
fixed.  These distinctions separate actuality from the two notions of
necessity studied in the cited systems
\cite{CrossleyHumberstone1977,DaviesHumberstone1980}.
\end{hbdefinition}

\begin{hbdefinition}{def:p08-character-content}{Context, character, and content}
In Kaplan's terminology, a sentence occurrence's \emph{character} is its rule
from contexts to contents, while a \emph{content} is what that rule supplies at
one context and is evaluated at a circumstance.  A context may contain an
agent, time, place, and world; only its world component is used in the present
comparison.  Thus, if context \(c\) designates \(y\), the content expressed by
formula \(\alpha\) at \(c\) is true at circumstance \(x\) exactly when
\(\langle W,V\rangle\models_x^y\alpha\).  Diagonal evaluation uses \(x=y\),
so it is true exactly when
\(\langle W,V\rangle\models_y^y\alpha\).  These are metalinguistic semantic
levels, not additional operators or newly named maps.  Character cannot replace
content because it still awaits a context, and content cannot replace character
because it no longer varies with context.  The distinction is used immediately
to state the two-dimensional payoff \cite{Kaplan1989}.
\end{hbdefinition}

\paragraph{Semantic payoff and epistemic limit.}
The clauses now make the relevant two-dimensional profile exact.  A sentence
may satisfy \(\langle W,V\rangle\models_u^u\alpha\) for every \(u\in W\),
while for some \(u,w\in W\) it fails off the diagonal:
\(\langle W,V\rangle\not\models_w^u\alpha\).  This is the
semantic pattern used in contingent-a-priori discussions: diagonal stability
does not entail necessity of the designated content.  The semantics alone does
not establish that a sentence is knowable a priori, and the Kaplanian
character/content distinction is not identified with either source operator
\(A\) or \(\mathcal F\).

\section{Quantified modal logic}

\begin{hbwarning}{def:p08-qml-family}{Quantified modal logic is a parameterized family}
There is no single calculus or semantics determined by the phrase
``quantified modal logic.''  A source presentation must specify its
propositional modal base, quantifier rules, domain variation, designation
policy, and local or global consequence.  The clauses below reproduce
Fitting's model and quantifier notation for the limited purpose of comparing
the actualist Barcan schemas; Garson is used only for the surrounding
comparison with other quantified-modal presentations.  No hybrid calculus or
system card is inferred from that semantic example
\cite{Fitting2002Alethic,FittingMendelsohn1998,Garson2001}.
\end{hbwarning}

\begin{hbdefinition}{def:p08-qml-model}{Fitting's varying-domain modal model}
In Fitting's notation a model is
\[
 \mathcal M=\langle\mathcal G,\mathcal R,\mathcal D,\mathcal I\rangle,
\]
where \(\mathcal G\) is a nonempty set of worlds,
\(\mathcal R\subseteq\mathcal G^2\), and
\(\mathcal D\) maps each \(\Gamma\in\mathcal G\) to a nonempty world-domain
\(\mathcal D(\Gamma)\).  The domain of the model is
\(\bigcup_{\Gamma\in\mathcal G}\mathcal D(\Gamma)\).  The interpretation
\(\mathcal I\) maps each predicate symbol \(P\) to its world-relative
extension \(\mathcal I(P)(\Gamma)\) over that model domain.  For Fitting's
special existence predicate \(E\),
\[
 \mathcal I(E)(\Gamma)=\mathcal D(\Gamma).
\]
Valuations \(v\) take variables into the model domain and do not vary with the
world.  For a formula \(\Phi\), truth is written
\(\mathcal M,\Gamma\forces_v\Phi\).  Ordinary
\((\forall x)\) and \((\exists x)\) are possibilist and range over the whole
model domain.  Fitting's actualist quantifiers are the exact abbreviations
\[
\begin{aligned}
 (\forall^{E}x)\Phi
 &\defeq (\forall x)[E(x)\supset\Phi],\\
 (\exists^{E}x)\Phi
 &\defeq (\exists x)[E(x)\land\Phi].
\end{aligned}
\]
Thus the relativized quantifiers range over \(\mathcal D(\Gamma)\) at world
\(\Gamma\), while modal satisfaction keeps \(v\) fixed when it moves along
\(\mathcal R\).  The model-domain convention is necessary to state
decreasing-domain failures with one world-independent valuation
\cite{Fitting2002Alethic,FittingMendelsohn1998}.
\end{hbdefinition}

\begin{hbdefinition}{def:p08-qml-domain-regimes}{Domain regimes and philosophical readings}
The domains are \emph{constant} when \(\mathcal D(\Gamma)\) has the same
value at every world; \emph{increasing} when
\(\Gamma\mathcal R\Delta\) implies
\(\mathcal D(\Gamma)\subseteq\mathcal D(\Delta)\);
\emph{decreasing} when \(\Gamma\mathcal R\Delta\) implies
\(\mathcal D(\Delta)\subseteq\mathcal D(\Gamma)\); and
\emph{arbitrarily varying} when no inclusion condition is imposed.  In this
section, \emph{actualist quantification} means Fitting's relativized
\(\forall^E,\exists^E\), while the primitive unrelativized quantifiers are
\emph{possibilist}; and the domain assignment is \emph{necessitist} exactly
when every world-domain is the whole model domain.  These are explicitly defined
interpretations of the semantic parameters, not consequences of the validity
table below.  The displayed BF/CBF table uses the actualist clauses.  Under
possibilist quantification over the model domain, varying the world-domains does not
generate that table; under necessitism, actualist and possibilist
quantification coincide with the constant-domain case.  This is the later use
that requires the three readings to remain distinct \cite{Williamson2013}.
\end{hbdefinition}

\begin{hbdefinition}{def:p08-barcan-schemas}{Barcan and converse Barcan schemas}
For a formula \(\Phi(x)\), the actualist \emph{Barcan formula}
\((\mathrm{BF}^{E})\) and actualist \emph{converse Barcan formula}
\((\mathrm{CBF}^{E})\) in Fitting's notation are respectively
\[
 (\forall^{E}x)\Box\Phi(x)\to\Box(\forall^{E}x)\Phi(x),
 \qquad
 \Box(\forall^{E}x)\Phi(x)\to(\forall^{E}x)\Box\Phi(x).
\]
The superscript records that this table uses the source's actualist
relativization.  With the primitive possibilist quantifier, local-domain
variation does not yield the table below.  The two names remain in this order;
their difference
is used in Table~\ref{tab:p08-barcan-domains} and
Figure~\ref{fig:p08-barcan-countermodel}; they are not interchanged when the
domain convention changes \cite{Barcan1946}.
\end{hbdefinition}

\begin{hbtheorem}{thm:p08-barcan-domain-validities}{cited=Fitting2002Alethic}{Actualist Barcan validity in Fitting's model}
Under Definitions~\ref{def:p08-qml-model}--\ref{def:p08-barcan-schemas}, the
validity pattern is exactly that displayed in
Table~\ref{tab:p08-barcan-domains}.  The decreasing-domain row uses the fixed
outer domain and rigid-assignment convention; without such a convention the
reversal is not a well-formed comparison.
\end{hbtheorem}

\begin{table}[!htbp]
\centering
\begin{tabular}{@{}lll@{}}
\toprule
World-domain regime & \(\mathrm{BF}^{E}\) & \(\mathrm{CBF}^{E}\) \\
\midrule
constant & valid & valid \\
increasing & may fail & valid \\
decreasing, with the fixed outer domain & valid & may fail \\
arbitrarily varying & may fail & may fail \\
\bottomrule
\end{tabular}
\caption{Actualist Barcan and converse-Barcan validity in Fitting's displayed
model notation (apparatus A-078).  ``May fail'' asserts the existence of a
model in the regime, not failure in every model.}
\label{tab:p08-barcan-domains}
\end{table}

\begin{hbexample}{def:p08-barcan-countermodel}{Increasing-domain failure of BF}
Let \(\mathcal G=\{\Gamma,\Delta\}\) and
\(\mathcal R=\{(\Gamma,\Delta)\}\).  Take
\(\mathcal D(\Gamma)=\{a\}\),
\(\mathcal D(\Delta)=\{a,b\}\),
\(\mathcal I(P)(\Gamma)=\varnothing\), and
\(\mathcal I(P)(\Delta)=\{a\}\).  Thus at \(\Delta\), \(P(a)\) is true and
\(P(b)\) is false.  At \(\Gamma\),
\((\forall^{E}x)\Box P(x)\) is true because the sole root-domain
object \(a\) satisfies \(P\) at the accessible world.  But
\(\Box(\forall^{E}x)P(x)\) is false because the new object
\(b\in\mathcal D(\Delta)\) fails \(P\) there.  Thus the model refutes
\(\mathrm{BF}^{E}\).  It does not refute \(\mathrm{CBF}^{E}\).
\end{hbexample}

\begin{figure}[!htbp]
\centering
\begin{tikzpicture}[>=Latex,node distance=4.6cm,font=\small]
  \node[draw,rounded corners,align=center,inner sep=6pt] (w)
    {\(\Gamma\)\\\(\mathcal D(\Gamma)=\{a\}\)\\
     \((\forall^{E}x)\Box P(x)\) true\\
     \(\Box(\forall^{E}x)P(x)\) false};
  \node[draw,rounded corners,align=center,inner sep=6pt,right=of w] (v)
    {\(\Delta\)\\\(\mathcal D(\Delta)=\{a,b\}\)\\
     \(P(a)\) true, \(P(b)\) false};
  \draw[->] (w) -- node[above]{\(\mathcal R\); domain increases by \(b\)} (v);
\end{tikzpicture}
\caption{An increasing-domain countermodel to \(\mathrm{BF}^{E}\) (apparatus A-142).  The two
displayed worlds and the one displayed edge are the whole frame.  The new
object at \(\Delta\) is outside the antecedent's root quantifier but inside the
consequent's successor quantifier.}
\label{fig:p08-barcan-countermodel}
\end{figure}

\paragraph{Source and navigation note.}
The normal-modal definitions, premise discipline for K, T, S4, and S5, bundled adequacy
statement, and three countermodel calculations migrate v1.1 Part X,
lines 13204--13339.  The rule schemas themselves occur only on the exact
registry sheets.  The canonical sketch is in Annex~\ref{annex:p8.modal.1}; the
new correspondence, filtration, bisimulation, neighborhood, and quantified
material is cited at its stated evidence ceiling.
